\documentclass[11pt]{article}
\usepackage[a4paper,margin=22mm]{geometry}
\usepackage{fontspec}
\setmainfont{DejaVu Serif}
\setsansfont{DejaVu Sans}
\usepackage{microtype}
\usepackage{xcolor}
\usepackage{hyperref}
\usepackage{enumitem}
\usepackage{parskip}
\definecolor{Accent}{HTML}{971D32}
\definecolor{Muted}{HTML}{666666}
\hypersetup{colorlinks=true,urlcolor=Accent,linkcolor=Accent}
\setlist{leftmargin=1.6em,itemsep=0.3em,topsep=0.35em}
\setcounter{tocdepth}{2}
\renewcommand{\contentsname}{Contents}
\newcommand{\sectionrule}{\par\vspace{-0.5em}\noindent\textcolor{Accent}{\rule{\linewidth}{0.6pt}}\par}
\begin{document}

\begin{center}
{\sffamily\bfseries\color{Accent} TOEFL WORD OF THE DAY\par}
\vspace{8mm}
{\fontsize{42}{48}\selectfont\bfseries contrite\par}
\vspace{2mm}
{\Large\sffamily\color{Accent} /kənˈtraɪt/\par}
\vspace{2mm}
{\small\sffamily Local audio phrase: contrite\par}
{\small\sffamily Audio file: audios/contrite.mp3\par}
\vspace{2mm}
{\large\sffamily\color{Muted} adjective\par}
\vspace{5mm}
{\sffamily sincerely sorry for wrongdoing\par}
\vspace{6mm}
{\large Contrite describes someone who feels and shows sincere guilt or regret for something wrong they have done.\par}
\vspace{6mm}
\begin{minipage}{0.84\textwidth}
\itshape “He appeared genuinely contrite after realizing how much his comments had hurt his colleagues.”
\end{minipage}\par
\vspace{10mm}
\textbf{Date:} August 21st 2026\quad\textbullet\quad \textbf{Level:} B1--mid-B2 TOEFL
\vfill
Created and compiled by \textbf{Amir Kooshky}\par
\vspace{2mm}
\href{https://t.me/KooshkyTOEFL}{Telegram Channel}\quad\textbullet\quad
\href{https://www.instagram.com/kooshkytoefl}{Instagram}\quad\textbullet\quad
\href{https://t.me/KooshkyTOEFL\_pv}{Personal Telegram}
\end{center}
\newpage
\tableofcontents
\newpage

\section{Meanings and usage}
\sectionrule
\subsection{Meaning 1: Sincerely sorry for wrongdoing}
\textbf{Part of speech:} adjective

\textbf{IPA:} /kənˈtraɪt/

\textbf{Pronunciation phrase:} contrite

\textbf{Local audio file:} audios/contrite.mp3

\textbf{Register:} formal; often used when discussing apologies, wrongdoing, responsibility, or remorse

\textbf{Definition:} feeling and showing sincere guilt or regret because you have done something wrong

\textbf{In simpler words:} truly sorry because you did something wrong

\textbf{Typical pattern:} be, seem, appear, or sound + contrite

\subsubsection{Natural examples}
\begin{enumerate}
  \item He appeared genuinely contrite after realizing how much his comments had hurt his colleagues.
  \item The company issued a contrite apology for failing to protect customers' personal information.
  \item She sounded contrite as she admitted that she had ignored the warning signs.
  \item The student was clearly contrite after being caught cheating on the examination.
  \item Although the manager apologized, employees did not believe that he was truly contrite.
  \item The politician adopted a contrite tone and accepted responsibility for the mistake.
  \item He returned the damaged equipment with a contrite expression and offered to pay for the repairs.
  \item The judge considered whether the defendant seemed genuinely contrite about his actions.
\end{enumerate}

\subsubsection{Nuance and implication}
\textbf{Central nuance:} Contrite is stronger than simply sorry. It suggests sincere regret together with recognition that you behaved wrongly.

\textbf{Usual implication:} A contrite person often admits responsibility and may try to repair the harm they caused.

\textbf{Compared with unrepentant:} A contrite person sincerely regrets wrongdoing. An unrepentant person does not show regret or believe that an apology is necessary.

\subsubsection{Grammar notes}
\textbf{How to use it:} Use it after a linking verb, as in She seemed contrite, or before a noun such as apology, tone, or expression.

\textbf{What can follow it:} It can be followed by about and the wrongdoing, although this pattern is less common than using contrite by itself.

\textbf{Common preposition:} When you need to name the wrongdoing, contrite about is possible, as in He was contrite about his behavior.

\textbf{Position in a sentence:} Words such as genuinely, deeply, truly, clearly, and suitably often appear before contrite.

\subsubsection{Useful patterns}
\textbf{Pattern:} be, seem, or appear contrite

\textbf{Use:} Use this when describing someone who shows sincere regret.

\textbf{Example:} The director appeared contrite during the public apology.

\textbf{Pattern:} sound contrite

\textbf{Use:} Use this when regret can be heard in someone's voice or words.

\textbf{Example:} She sounded contrite when she admitted that the mistake was hers.

\textbf{Pattern:} a contrite apology

\textbf{Use:} Use this for an apology that strongly expresses guilt and regret.

\textbf{Example:} The newspaper published a contrite apology for the inaccurate report.

\textbf{Pattern:} a contrite tone or expression

\textbf{Use:} Use this for a voice or facial expression that shows sincere regret.

\textbf{Example:} He spoke in a contrite tone and promised not to repeat the mistake.

\textbf{Pattern:} genuinely or deeply contrite

\textbf{Use:} Use these words to emphasize that the regret is sincere or strong.

\textbf{Example:} She seemed genuinely contrite about the harm she had caused.

\subsubsection{Common wrong usage}
\paragraph{Correction 1}
\textbf{Wrong:} She was contrite because she was disappointed that the concert was canceled.

\textbf{Correct:} She was disappointed because the concert was canceled.

\textbf{Why:} Contrite is used for regret about your own wrongdoing, not ordinary disappointment about an unfortunate event.

\paragraph{Correction 2}
\textbf{Wrong:} He gave a contritely apology.

\textbf{Correct:} He gave a contrite apology.

\textbf{Why:} Use the adjective contrite before the noun apology.

\paragraph{Correction 3}
\textbf{Wrong:} The employee contrited after making the mistake.

\textbf{Correct:} The employee became contrite after making the mistake.

\textbf{Why:} Contrite is an adjective, not a verb.

\paragraph{Correction 4}
\textbf{Wrong:} She was contrite of her behavior.

\textbf{Correct:} She was contrite about her behavior.

\textbf{Why:} If you name the wrongdoing after contrite, use about rather than of.

\section{Word family}
\sectionrule
\subsection{contrition}
\textbf{Part of speech:} uncountable noun

\textbf{IPA:} /kənˈtrɪʃən/

\textbf{Definition:} sincere guilt and regret for something wrong that you have done

\textbf{Relationship to the headword:} This noun names the feeling of sincere guilt or regret described by contrite.

\textbf{Grammar pattern:} show, express, or feel + contrition, or contrition for + wrongdoing

\textbf{Usage note:} It is formal and normally used without a or an.

\subsubsection{Examples}
\begin{enumerate}
  \item He showed genuine contrition for the harm his actions had caused.
  \item The apology expressed little real contrition.
  \item Her willingness to accept responsibility suggested genuine contrition.
\end{enumerate}

\subsection{contritely}
\textbf{Part of speech:} adverb

\textbf{IPA:} /kənˈtraɪtli/

\textbf{Definition:} in a way that shows sincere guilt or regret for wrongdoing

\textbf{Relationship to the headword:} This adverb describes speaking or behaving in a contrite way.

\textbf{Grammar pattern:} contritely + admit, apologize, say, or acknowledge

\textbf{Usage note:} It is less common than contrite and contrition.

\subsubsection{Examples}
\begin{enumerate}
  \item He contritely admitted that he had ignored the safety procedure.
  \item She apologized contritely and offered to correct the mistake.
\end{enumerate}

\section{Common collocations}
\sectionrule
\subsection{genuinely contrite}
\textbf{Applies to:} Sincerely sorry for wrongdoing

\textbf{Meaning:} sincerely and honestly sorry for wrongdoing

\textbf{Example:} The committee believed that the employee was genuinely contrite.

\subsection{deeply contrite}
\textbf{Applies to:} Sincerely sorry for wrongdoing

\textbf{Meaning:} feeling very strong guilt or regret

\textbf{Example:} He was deeply contrite about the damage his decision had caused.

\subsection{appear contrite}
\textbf{Applies to:} Sincerely sorry for wrongdoing

\textbf{Meaning:} seem to feel sincere regret

\textbf{Example:} The executive appeared contrite during the press conference.

\subsection{seem contrite}
\textbf{Applies to:} Sincerely sorry for wrongdoing

\textbf{Meaning:} give the impression of feeling sincere regret

\textbf{Example:} She seemed contrite after learning how seriously her actions had affected others.

\subsection{sound contrite}
\textbf{Applies to:} Sincerely sorry for wrongdoing

\textbf{Meaning:} speak in a way that shows sincere regret

\textbf{Example:} The manager sounded contrite when he accepted responsibility for the failure.

\subsection{contrite apology}
\textbf{Applies to:} Sincerely sorry for wrongdoing

\textbf{Meaning:} an apology expressing sincere guilt and regret

\textbf{Example:} The organization issued a contrite apology to the affected families.

\subsection{contrite tone}
\textbf{Applies to:} Sincerely sorry for wrongdoing

\textbf{Meaning:} a way of speaking that shows sincere regret

\textbf{Example:} He adopted a contrite tone when discussing the mistake.

\subsection{contrite expression}
\textbf{Applies to:} Sincerely sorry for wrongdoing

\textbf{Meaning:} a facial expression showing guilt or regret

\textbf{Example:} Her contrite expression suggested that she understood the seriousness of the situation.

\subsection{show contrition}
\textbf{Applies to:} Sincerely sorry for wrongdoing

\textbf{Meaning:} demonstrate sincere regret for wrongdoing

\textbf{Pattern:} show contrition for + wrongdoing

\textbf{Example:} The official showed genuine contrition and accepted full responsibility.

\textbf{Usage note:} Contrition is the noun form.

\subsection{express contrition}
\textbf{Applies to:} Sincerely sorry for wrongdoing

\textbf{Meaning:} state or demonstrate sincere regret

\textbf{Pattern:} express contrition for + wrongdoing

\textbf{Example:} The company expressed contrition for its failure to follow safety regulations.

\textbf{Usage note:} Contrition is the noun form.

\section{Synonyms and antonyms}
\sectionrule
\subsection{Close synonyms}
\subsubsection{remorseful}
\textbf{Part of speech:} adjective

\textbf{Applies to:} Sincerely sorry for wrongdoing

\textbf{Shared meaning:} Both describe sincere regret and guilt about wrongdoing.

\textbf{Difference:} Remorseful is very close in meaning and emphasizes painful guilt about something wrong you have done. Contrite is more formal and often suggests that the regret is openly shown through an apology or behavior.

\textbf{Example:} The driver was deeply remorseful after causing the accident.

\subsubsection{repentant}
\textbf{Part of speech:} adjective

\textbf{Applies to:} Sincerely sorry for wrongdoing

\textbf{Shared meaning:} Both describe someone who recognizes and regrets wrongdoing.

\textbf{Difference:} Repentant emphasizes regretting wrongdoing and wanting to change one's behavior. It is also common in religious contexts. Contrite emphasizes sincere guilt and humility after doing wrong.

\textbf{Example:} The repentant offender promised to change his behavior.

\subsubsection{apologetic}
\textbf{Part of speech:} adjective

\textbf{Applies to:} Sincerely sorry for wrongdoing

\textbf{Shared meaning:} Both can describe someone who shows regret.

\textbf{Difference:} Apologetic means showing that you are sorry, often for a mistake or inconvenience. Contrite is stronger and suggests deeper guilt or moral regret.

\textbf{Example:} The receptionist was apologetic about the long delay.

\subsection{Direct or useful antonyms}
\subsubsection{unrepentant}
\textbf{Part of speech:} adjective

\textbf{Applies to:} Sincerely sorry for wrongdoing

\textbf{Opposition:} An unrepentant person shows no regret for wrongdoing, while a contrite person sincerely regrets what they have done.

\textbf{Example:} He remained unrepentant and insisted that he had done nothing wrong.

\subsubsection{remorseless}
\textbf{Part of speech:} adjective

\textbf{Applies to:} Sincerely sorry for wrongdoing

\textbf{Opposition:} A remorseless person feels or shows no guilt, while a contrite person feels sincere guilt and regret.

\textbf{Example:} The character is portrayed as remorseless despite the harm he causes.

\section{Words learners may confuse}
\sectionrule
\subsection{apologetic}
\textbf{Part of speech:} adjective

\textbf{Why learners confuse them:} Both describe someone who expresses regret or says they are sorry.

\textbf{Key difference:} Apologetic means showing that you are sorry and can be used for relatively small mistakes or inconveniences. Contrite usually suggests deeper guilt about genuine wrongdoing.

\textbf{contrite:} He was contrite after admitting that he had deliberately lied.

\textbf{apologetic:} The waiter was apologetic about bringing the wrong order.

\subsection{contrition}
\textbf{Part of speech:} adjective versus uncountable noun

\textbf{Why learners confuse them:} The words belong to the same family and express the same basic idea in different grammatical forms.

\textbf{Key difference:} Contrite is an adjective describing a person, expression, tone, or apology. Contrition is the noun for the feeling of sincere guilt and regret.

\textbf{contrite:} The manager seemed genuinely contrite.

\textbf{contrition:} The manager expressed genuine contrition.

\section{Etymology}
\sectionrule
\textbf{Origin:} Contrite comes from Latin contritus, meaning worn down or crushed.

\textbf{Development:} The physical idea of being crushed developed into the figurative idea of being humbled or emotionally affected by guilt for wrongdoing.

\textbf{Memory link:} A contrite person can be imagined as emotionally humbled or crushed by regret.

\section{Practice exercises}
\sectionrule
\subsection{Word family choice}
Choose the best member of the word family for each blank.

\begin{enumerate}
  \item The official expressed genuine \_\_\_\_\_ for misleading the public.
  \item She \_\_\_\_\_ admitted that the mistake was entirely her responsibility.
  \item The employee appeared genuinely \_\_\_\_\_ after the incident.
\end{enumerate}

\subsection{Correct or incorrect?}
Decide whether each sentence is correct. Correct the inaccurate ones.

\begin{enumerate}
  \item He seemed genuinely contrite after admitting that he had lied.
  \item She felt contrite because the weather ruined her vacation.
  \item The company issued a contrite apology for the data breach.
  \item The manager gave a contritely response.
  \item The official showed genuine contrition for his actions.
\end{enumerate}

\subsection{Collocation practice}
Complete each sentence with a natural collocation.

\begin{enumerate}
  \item The organization issued a contrite \_\_\_\_\_ for failing to protect customer data.
  \item The employee seemed genuinely \_\_\_\_\_ after realizing the harm he had caused.
  \item The official expressed contrition \_\_\_\_\_ misleading the public.
  \item She spoke in a contrite \_\_\_\_\_ and accepted responsibility for the mistake.
\end{enumerate}

\subsection{Complete answer key}
\subsubsection{Word family choice}
\begin{enumerate}
  \item \textbf{contrition.} The official expressed genuine contrition for misleading the public. \emph{Use the noun contrition after expressed.}
  \item \textbf{contritely.} She contritely admitted that the mistake was entirely her responsibility. \emph{Use the adverb contritely to describe how she admitted the mistake.}
  \item \textbf{contrite.} The employee appeared genuinely contrite after the incident. \emph{Use the adjective contrite after appeared.}
\end{enumerate}

\subsubsection{Correct or incorrect?}
\begin{enumerate}
  \item \textbf{Correct} \emph{Contrite correctly describes sincere regret for wrongdoing.}
  \item \textbf{Incorrect}. She felt disappointed because the weather ruined her vacation. \emph{Contrite describes guilt about your own wrongdoing, not sadness about something unfortunate.}
  \item \textbf{Correct} \emph{Contrite apology is a natural expression for an apology showing sincere regret.}
  \item \textbf{Incorrect}. The manager gave a contrite response. \emph{Use the adjective contrite before the noun response.}
  \item \textbf{Correct} \emph{Show contrition for something is a natural formal pattern.}
\end{enumerate}

\subsubsection{Collocation practice}
\begin{enumerate}
  \item \textbf{apology.} The organization issued a contrite apology for failing to protect customer data. \emph{A contrite apology expresses sincere regret for wrongdoing.}
  \item \textbf{contrite.} The employee seemed genuinely contrite after realizing the harm he had caused. \emph{Genuinely contrite means sincerely sorry for wrongdoing.}
  \item \textbf{for.} The official expressed contrition for misleading the public. \emph{Use contrition for followed by the wrongdoing.}
  \item \textbf{tone.} She spoke in a contrite tone and accepted responsibility for the mistake. \emph{A contrite tone is a way of speaking that shows sincere guilt or regret.}
\end{enumerate}

\vfill
\begin{center}
\small\color{Muted} Created and compiled by \textbf{Amir Kooshky}\\
\href{https://t.me/KooshkyTOEFL}{Telegram Channel}\quad\textbullet\quad
\href{https://www.instagram.com/kooshkytoefl}{Instagram}\quad\textbullet\quad
\href{https://t.me/KooshkyTOEFL\_pv}{Personal Telegram}
\end{center}
\end{document}
